\documentclass[12pt]{article}
\usepackage[utf8]{inputenc}
\usepackage[T1]{fontenc}
\usepackage{hyperref}
\title{Yalarnnga Language Resources for K-12 Curriculum Development}
\author{Compiled by Claude Sonnet 5.5 (claude-sonnet-5-5) for Andr\'as Kornai}
\date{}
\begin{document}
\maketitle

\section{Introduction}
Yalarnnga (ISO 639-3 ylr; Glottocode yala1262) is an extinct Pama-Nyungan language of north-western Queensland, Australia. Search results give 1980 as the year the last native speaker died. Any curriculum would therefore be a revitalization or heritage curriculum built from documentary sources. I found no school programme, community-run course or teaching materials in my searches.

\section{Primary Educational Resources}
I found no primers, textbooks or classroom materials for Yalarnnga. Any teaching materials would have to be developed from the reference works below, ideally together with descendant communities and Indigenous language bodies.

\section{Linguistic and Reference Materials}
\begin{itemize}
\item Breen, Gavan and Blake, Barry. \emph{The grammar of Yalarnnga: a language of western Queensland} (Pacific Linguistics, 2007). The ANU Open Research Repository offers it as a free PDF. I did not open the PDF, so I cannot describe its contents beyond the title and the search summary. \url{https://openresearch-repository.anu.edu.au/bitstream/1885/146740/1/584_Breen&amp;Blake.pdf}
\item The ASJP database has a Yalarnnga wordlist. It cites Breen's vocabulary and Blake's 1971 morphology outline. \url{https://asjp.clld.org/languages/YALARNNGA}
\item The Grambank entry for Yalarnnga gives grammatical feature data. \url{https://grambank.clld.org/languages/yala1262}
\item The English Wikipedia article summarizes the language and its sources, including Eglinton's 1886 account in Curr's \emph{The Australian Race}. I did not verify those sources against the originals. \url{https://en.wikipedia.org/wiki/Yalarnnga_language}
\end{itemize}

\section{Religious and Community Resources}
The searches turned up no scripture translations and no community organisation devoted to Yalarnnga.

\section{Digital and Audio Resources}
I found no audio or video recordings and no apps. The ASJP wordlist and the PDF grammar above are the only digital materials I found.

\section{Curriculum Development Context}
Material is very limited: one grammar, wordlists, and short historical records. Any K-12 use would be heritage awareness, with words and simple phrases drawn from the Breen and Blake grammar. It would need consultation with descendant communities and Queensland Indigenous language bodies. I did not search for such bodies, so none is named here.
\end{document}